%% LyX 2.4.5 created this file.  For more info, see https://www.lyx.org/.
%% Do not edit unless you really know what you are doing.
\documentclass[11pt,english]{article}
\usepackage[T1]{fontenc}
\usepackage[utf8]{inputenc}
\usepackage{color}
\usepackage{microtype}
\usepackage{xurl}
\usepackage{array}
\usepackage{float}
\usepackage{booktabs}
\usepackage{amsmath}
\usepackage{graphicx}
\usepackage{geometry}
\geometry{verbose,tmargin=1.1in,bmargin=1.1in,lmargin=1.1in,rmargin=1.1in}
\usepackage{setspace}
\onehalfspacing
\setlength{\emergencystretch}{3em}

\makeatletter

%%%%%%%%%%%%%%%%%%%%%%%%%%%%%% LyX specific LaTeX commands.
%% Because html converters don't know tabularnewline
\providecommand{\tabularnewline}{\\}
%% A simple dot to overcome graphicx limitations
\newcommand{\lyxdot}{.}


%%%%%%%%%%%%%%%%%%%%%%%%%%%%%% User specified LaTeX commands.
%
\usepackage{amsfonts}\newtheorem{prop}{Proposition}
\newtheorem{cor}{Corollary}
\newtheorem{lem}{Lemma}
\interfootnotelinepenalty=10000
\setcounter{MaxMatrixCols}{30}
\providecommand{\U}[1]{\protect\rule{.1in}{.1in}}
\newcommand{\rev}[1]{\textcolor{blue}{#1}}
\newcommand{\revmath}[1]{{\color{blue}#1}}
\newcommand{\figredo}[1]{\textcolor{blue}{[Figure to redo: #1]}}
\newenvironment{proof}[1][Proof]{\noindent\textbf{#1.} }{\ \rule{0.5em}{0.5em}}

\newtheorem{theorem}{Theorem}
\newtheorem{acknowledgement}{Acknowledgement}
\newtheorem{algorithm}[theorem]{Algorithm}
\newtheorem{axiom}[theorem]{Axiom}
\newtheorem{assumption}{Assumption}
\newtheorem{case}[theorem]{Case}
\newtheorem{claim}[theorem]{Claim}
\newtheorem{conclusion}[theorem]{Conclusion}
\newtheorem{condition}[theorem]{Condition}
\newtheorem{conjecture}[theorem]{Conjecture}
\newtheorem{corollary}[theorem]{Corollary}
\newtheorem{criterion}[theorem]{Criterion}
\newtheorem{definition}{Definition}
\newtheorem{example}[theorem]{Example}
\newtheorem{exercise}[theorem]{Exercise}
\newtheorem{lemma}{Lemma}
\newtheorem{notation}[theorem]{Notation}
\newtheorem{problem}[theorem]{Problem}
\newtheorem{proposition}{Proposition}
\newtheorem{remark}[theorem]{Remark}
\newtheorem{solution}[theorem]{Solution}
\newtheorem{summary}[theorem]{Summary}
\usepackage{lscape}

\makeatother

\usepackage{babel}
\begin{document}
\title{NIMBYism and the Housing Supply\thanks{We thank Stephane Bonhomme, Fatih Guve\textcolor{black}{nen, Serdar
Ozkan and Roman Sustek for help for their comments as well as with
data and calibration issues. We also thank participants in seminars
at Durham University, Monash University, University of Sydney, CEMAP,
WAMS, OzMac, SCE a}nd SEHO. All errors remaining are our own.}}
\author{Isaac Gross\thanks{isaac.gross@monash.edu}\\
{\small Monash University} \and David Chivers\thanks{david.chivers@durham.ac.uk\textbf{ }}
\\
{\small Durham University} \\
}
\maketitle
\begin{abstract}
\ We develop a political economy model of the housing market in which
the supply of new homes is shaped by a democratic vote among households.
Voter support for additional housing is determined by the effect this
vote has on households' expected discounted utility via changes in
house prices. Households that own housing, and therefore typically
experience financial losses from additional supply, generally oppose
new housing: the 'Not In My Backyard' (NIMBY) effect. Households that
wish to buy a house and would benefit from cheaper housing support
more supply. We find that support for additional supply is highest
among households that are younger, low-income or are currently renting.
\rev{We first characterize a sequence of stationary political-economic
equilibria indexed by the age distribution of households. We then
study the transition generated by a temporary baby boom under political
rational expectations, so that the future path of house prices and
future political support are solved jointly with household decisions.
In that environment, a large cohort initially supports more housing
when young and then supports tighter supply as it ages into homeownership,
redistributing housing access and wealth across generations. Calibrating
the model to US demographic data, we find that demographic change
exerts substantial upward pressure on house prices, and the stationary
benchmark implies further pressure from population aging in the decades
ahead.}

\

\textbf{JEL Classification}: E24, J10, R31.

\textbf{Keywords}: Housing supply, inequality, political economy
\end{abstract}
\pagebreak{}

\section{Introduction}

The ``Not In My Backyard'' (NIMBY) effect describes the phenomenon
of local residents opposing new housing developments in their neighborhood.
As the regulation of housing supply is often decentralized to local
governments, residents can exert significant political power over
any new developments within their area. The aggregate effect of NIMBYism
can thus negatively impact the housing supply within an economy. In
the US, for example, land is divided into ``zones'' which determine
the size and use of new developments.\footnote{Zoning laws were originally developed in the 1900s in order to improve
``light and air'' in residential areas - see Fischer (2002) for
a historical overview. Land use regulation has also been used by white
residents in the US to maintain and increase racial segregation (Sahn
2021, Trounstine 2020).} Only detached, single-family homes are allowed to be built in ``R1''
residential zones, even if high demand for multi-home housing (such
as apartments) exists. Those who own homes within R1 residential areas
have strong incentives to maintain the status quo at the expense of
non-homeowners, who would benefit from multi-home housing in order
to get on the housing ladder (Manville et al. 2020).

We develop a political economy model of the housing market in order
to estimate the impact of NIMBYism and show how changing demographic
trends can explain the recent rise in house prices. Specifically,
we allow housing supply in a standard life cycle model to be determined
by a democratic process at the local level: households vote according
to their personal preferences and the final outcome is determined
by the median voter. \rev{Our quantitative analysis has two parts.
First, we solve a sequence of stationary political-economic equilibria
indexed by the demographic composition of the population. This benchmark
isolates how the age distribution shifts the median voter's preferred
level of housing supply and therefore the equilibrium house price.
Second, we study the transition generated by a temporary baby boom
under political rational expectations, so that households' continuation
values incorporate the future path of house prices and future political
support implied by the model itself. This lets us distinguish the
long-run demographic benchmark from the fully dynamic redistribution
created by a large cohort as it moves from renting into homeownership
and then into net seller status.}

Our theoretical model follows the standard life cycle housing model
as described by Iacaviello (2005). Over the life-cycle, an individual
household typically rents housing until they have saved enough to
meet their borrowing constraint and be able to put a deposit down
on a house. Households also participate in the housing market by either
upsizing (as their wealth grows or positive income shocks occur) or
selling (if they receive a negative income shock or die). Instead
of fixing the housing supply (as in Iacaviello 2005) or having it
produced inelastically by a construction sector (as in Kaplan et al.
2020), we allow the supply of housing to be determined by a democratic
vote of local households.

Household preferences over the supply of housing are formed according
to the impact that a small change in price will have on their expected
discounted utility. Households who gain from marginally higher prices
will vote against new developments, whereas those who marginally lose
out from higher house prices will vote in favor of new developments.\footnote{Households assume that increasing the housing supply decreases house
prices within their local area. Although increasing housing supply
may increase house prices via agglomeration effects, these are dispersed
across a larger area (Duranton and Kerr, 2015) and so will be ignored
by households voting within decentralised local government authorities
that make up a greater metropolitan area. This dynamic underpins the
NIMBY position: in principle, they may support increased housing supply
- just not in their own backyard.. } NIMBYist views in our enviroment are thus largely determined by households'
ownership of housing and their expected probability of being a net-buyer
in future periods. NIMBYist households are disproportionately older,
wealthier, have higher incomes and are more likely to be homeowners,
which is in line with the empirical literature surrounding NIMBYism
(see Whittemore \& BenDor (2018), Scally and Tighe (2015), and Hankinson
(2018) for a range of surveys on the subject).

Our model primarily focuses on the financial incentives of NIMBYism,
as households' own homes are the primary asset for most households.\footnote{Excluding the top 1\% of the wealth distribution, home equity makes
up 28.9\% of total US household wealth and is comparable only to retirement
accounts 32.8\%. The next largest asset, stocks and mutual funds,
makes up only 10.2\% of total household wealth (see Eggleston et al.
2020).} However there are several non-financial motivations that may lead
a household to oppose new developments - ranging from infrastructure
capacity concerns to simply disliking the proposed style of architecture
or neighborhood characteristics (Ahlfeldt 2011, Bosetti and Sims 2016)\textit{.}
We allow for these non-financial preferences to affect households'
votes but assume these preferences are stable over the life cycle\textit{.}
The financial inventive for NIMBYism, however, does vary over the
life cycle and is analogous to an insider-outsider problem: outsiders
(renters) would prefer more houses to be built until they become an
insider (homeowner), and then seek to restrict supply.

As homeownership rates and thus support for new developments changes
over the life cycle, the demographic make-up of an economy will play
an important role in the aggregate effect of NIMBYism. During the
mid 1940s, a baby boom occurred in a number of OECD countries including
the US, UK, New Zealand, Australia, France and Canada (Greenwood et
al. 2005). This baby boom lasted until the subsequent baby bust of
the mid 1960s when birth rates returned to trend. The effect of the
spike in the birth rate led to a generation known as the Baby Boomers
which has taken up a significant share of the population over time.
As those who are close in age will have similar preferences over housing
on average (as a result of being at a similar stage of the life-cycle),
those born during the baby boom will have outsized political power
and thus have a greater influence over the aggregate housing stock.
This is perhaps one reason why rents are increasing relative to median
income over the long run as shown in Figure \ref{Figure 1}.
\begin{figure}[H]
\noindent\begin{raggedright}
\caption{\textit{\emph{Rent to Median Income Ratio Over Time}}}
\par\end{raggedright}
\begin{centering}
\emph{\includegraphics[scale=0.4]{\string"Figures/1. Rent to Median Income\string".eps}}
\par\end{centering}
\begin{centering}
{\tiny Data Source: CPS (2022)}{\tiny\par}
\par\end{centering}
\raggedright{}\label{Figure 1}
\end{figure}
\emph{ }It is also why the US homeownership rate has been falling
for younger people over time as shown in figure \ref{Figure 2}.\emph{
}
\begin{figure}[H]
\caption{Home Ownership Rate for Different Age Cohorts}

\begin{centering}
\emph{\includegraphics[scale=0.4]{\string"Figures/2. Home Ownership Rate\string".eps}}
\par\end{centering}
\begin{centering}
{\tiny Data Source: CPS (2022)}{\tiny\par}
\par\end{centering}
\label{Figure 2}
\end{figure}

In this paper, we estimate how the changing demographic shares across
aging economies have changed the political-economic equilibrium in
the housing market. As economies age, the balance of political power
shifts in favor of NIMBYist households resulting in lower housing
supply and higher house prices. We find that house prices increased
by around 15 per cent due to these political economic effects since
the 1950s. Using demographic projections, as shown in figure \ref{Figure 3}
we find that the aging economy is expected to keep the house price
at or above it's current level until 2100.

Finally, we use our political economic model of the housing market
to examine the distributional consequences of recent demographic shocks.
We find that a temporary boom in the birthrate will initially decrease
inter and intra-generational wealth inequality as the boom generation
initially supports lower house prices when they are disproportionately
young renters. As the cohort ages, this effect reverses and this cohort
greater political pressure against the supply of new housing. This
decrease in the support for housing supply, and thus higher house
price, leads to a persistent scarring effect on future generations
who are unable to achieve home ownership at the same rate as previous
generations, even after the initial baby boom generation has completely
died off.
\begin{figure}[H]
\caption{\textit{\emph{Age Share of Population Over Time}}}

\begin{centering}
\includegraphics[scale=0.5]{\string"Figures/3. Projections by age group, median immigration\string".eps}
\par\end{centering}
\centering{}{\tiny Data Source: UN (2019). Dashed line indicates population
forecasts for year 2020 onward assuming mid level immigration} \label{Figure 3}
\end{figure}


\section{Related Literature }

This paper is related to the following strands of the literature.
Firstly, there is growing literature surrounding the impact of regulation
on the housing supply. The concept of NIMBYism is often attributed
to Frieden (1979), who noted a rise in environmental regulations to
impede new housing developments after the 1970s. Gyourko et al. (2013)
support this view and argue that before the 1970s, housing construction
tended to follow areas of high price growth. However, establishing
a direct link between regulation and the housing supply is difficult
due to the idiosyncratic forms that regulations can take in different
localities (Gyourko and Malloy 2015). To get round this, Glaeser et
al. (2006) offer indirect evidence of excess regulation by noting
the existence of a relatively large wedge between prices and marginal
constructions costs. The difficulty in providing direct evidence stems
from the endogeneity between land use regulation and the housing market.
Recent papers have used boundary discontinuity design to address this
issue (Kahn et al. 2010, Turner et al. 2014, Kulka 2019). For example,
Anagol et al. (2021) exploit a zoning reform in Sao Paulo and show
that housing developers apply for more permits in blocks with higher
maximum densities. In their equilibrium model of housing supply this
leads to a 2.2\% increase in the total housing stock and 0.5\% average
price reduction. Similarly, Song (2021) constructs a nationwide data
set of minimum lot area estimates and finds that a doubling the US
minimum lot area size would increase house prices by 14\% and rents
by 6\%.

In theory, an increase in the supply of housing units may increase
house prices via improved amminties and neighbourhood qualility, or
via agglomaration effects which increase economic activity (Herkenhoff
et al. 2018, Chatterjee and Eyigungor 2017, Titman and Guozhong 2024).
In the immediate local vicinity, however, these demand effects do
not appear to dominate the supply channel. For example, Li (2002)
finds that within 500 ft, for every 10\% increase in the housing stock,
rents decrease by 1\% and for every 10\% increase in the condo stock,
condo sales prices decrease by 0.9\%. Similarly, Pennington (2021)
finds that rents fall by about 2\% for properties within 100 meters
of a new housing project, and the effects decrease to zero when 1.5
km away. This highlights a major feature of NIMBYism in that it is
ultimately a coordination problem: the impact of extra housing within
a homeowners immediate vicinity will negatively affect prices, whereas
the benefits of agglomation within the wider area can only occur if
local residents accept extra housing in their back yard. As a result
an emerging political economy literature has started to address this
issue (see for example, Parkhomenko 2023, Bunten 2017, Duranton and
Puga 2022 and Albouy et al., 2019). Using a difference in difference
approach, Mast (2024) shows that when elections of town councils were
localised from \textquotedblleft at-large\textquoteright\textquoteright{}
to \textquotedblleft ward\textquotedblright{} areas, housing permits
decreased by 20\%.

This paper is also related to an older strand of literature on political
redistribution, such as Alesina and Rodrik (1994) and Persson and
Tabellini (1994). The main idea of these papers is that when income
and resources are distributed heterogeneously, the majority may vote
for redistribute policies. These papers rely on a voting mechanism,
the median-voter theorem, to determine what policy will be pursued.
The theorem shows that when voting takes place over a single issue,
and preferences are singled peaked, then the policy chosen under a
democracy will be the one chosen by the median-voter.

Housing politics is a good candidate for the median-voter theorem
as voting take place at local level and usually on specific propositions
related to relaxing zoning regulations. The idea of homeowners as
specific voting class in local government decisions as a result of
financial interests is related to the ``homevoter hypothesis'' (Fischel
1999, 2001). Einstein et al. (2019) examine planning meetings in Massachusetts
and find that homeowners and older residents are overly representative
in planning meetings. This is in line with voter turnout in local
elections which in general tend to be from these two groups (Kogan
et al. 2018, Oliver and Ha 2007). Similarly, Jiang (2018) finds that
voter turnout increase for mayoral elections, relative to presidential
elections, when home-ownership rates increase. In addition, Cheung
and Meltzer (2013) find that the prevalence of Homeowners Associations
is positively associated with an increase in land use regulations.
Similarly, Hilber and Robert-Nicoud (2013) find a positive correlation
between the degree of physical development in an area and land use
regulation. Both Ahlfeldt and Maennig (2015) as well as Dehring and
Ward (2008) use public referendums to directly test the homevoter
hypothesis. Ahlfeldt and Maennig (2015) use a referendum on the expansion
of an airport in Berlin as a natural experiment and find greater support
among homeowners as opposed to leaseholders. The support for the proposal
comes from two effects: a consumption effects (though use of the amenity)
as well as a wealth effect (via house price increases). Dehring and
Ward (2008) use pre-vote changes in residential property values to
test the probability of support for the new stadium for the Dallas
Cowboys and find between a 0.9\% and 1.2\% increase in support for
every \$1000 dollar increase in house prices.

Numerous surveys of voters have attempted to describe the demographics
of voters who oppose new housing supply. Whittemore and BenDor (2019)
find that older, less educated were more likely to oppose new developments,
though they found no significant relationship between home ownership
and NIMBYist views. By contrast Hankinson (2018) and Scally and Tighe
(2015) found that home owners were significantly more likely to oppose
new housing supply - particularly in their locality.

Finally, this paper is related to the literature surrounding life-cycle
models with housing markets. Generally, life-cycle models assume exogenous
house prices due to computational complexity (Fernandez-Villaverde
and Krueger (2011), Iacoviello and Pavan (2013), Yang (2009)). There
are, however, relatively few papers where aggregate shocks affect
house prices (Justiniano et al. 2015). Our paper is most closely related
to Ortalo-Magné and Prat (2014), who use a median-voter theorem in
an overlapping\textendash generations economy to show that increasing
homeownership rates decrease housing affordability. Similarly, Mankiw
and Weil (1989) predict the growth rate of housing demand via changes
in the age distribution alone and more recently Mumtaz and ¦ustek
(2023) find similar results except housing demand peaks in the 2000s.\textcolor{red}{{} }

\section{The Model}

In this section we outline a life-cycle model with housing and debt
choices. The model features heterogeneous, life-cycle households that
optimize their consumption of non-durable consumption, housing and
saving in liquid assets. A number of factors influence households
preferences on house prices and thus their political support for additional
housing supply. These include household income and wealth as well
as portfolio allocation and age over the life cycle. We use the model
to study the housing market response to an exogenous change in demographics.
Following the shock, the behavior of owner-occupiers and renters determines
equilibrium house prices.

A consumer of age $j$ enters the economy at age 25 and lives for
further $J$ periods and maximizes lifetime expected utility which
consists of consumption $c_{t}$, housing services, $h_{t}$ and a
bequest $v\left(w_{J+1}\right)$. Housing services depends on whether
they own or rent the property. Individuals also have access to a liquid
asset that pays an interest $r_{t+1}$. The household's lifetime utility
function is given by

\begin{equation}
\sum^{J}_{j=0}\beta^{j}u(c_{j},h_{j})+\beta^{J+1}v\left(w_{J+1}\right)
\end{equation}
where households optimize over consumption, $c_{t}$, liquid assets,
$a_{t}$, and housing, $h_{t}$ at each age $j$. The flow utility
is increasing and strictly concave in $c$ and $h$. Household receive
utility from wealth they bequest $w_{J+1}$. We assume the functional
form

\begin{equation}
u(c_{t},h_{t})=\frac{\left(c^{1-\zeta}_{t}h^{\zeta}_{t}\right)^{1-\gamma}}{1-\gamma}
\end{equation}
where the consumption of housing stock consists of either the amount
of housing owned or rented

\begin{equation}
h_{t}=\begin{cases}
\theta^{rent}h^{rent}_{t} & \mbox{if }h^{own}_{t}=0\\
h^{own}_{t} & \mbox{if }h^{own}_{t}>0.
\end{cases}
\end{equation}
The curvature parameter $\gamma$ determines households' risk aversion
and inter-temporal elasticity of substitution and $\zeta$ denotes
housing's share of total consumption. Rental housing provides a stream
of utility that is discounted by $\theta^{rent}$ compared to self-owned
housing. Finally there is a warm glow bequest function when households
die that provides utility
\begin{equation}
v(w)=\psi\frac{\left(w+\overline{w}\right)^{1-\gamma}}{1-\gamma}
\end{equation}
which generates a bequest function with decreasing marginal utility
with respect to wealth, $w$, governed by the parameters $\overline{w}$
and $\psi$. Households supply one unit of labor inelastically and
are paid a wage $y_{t,j}$. The wage they receive varies according
to their productivity which is driven by the exogenous process

\begin{eqnarray}
\text{ln}y_{t,j} & = & g_{j}+z_{t}\text{ for }j\leq J^{retired}
\end{eqnarray}
where $z_{t}$ is a Markov process calibrated to yield realistic inequality
(as per Kindermann and Krueger (2022)) and $g_{j}$ is a deterministic
age profile. The stochastic process comprises of an AR(1) process
for lowest income states, with two additional top-income states which
incorporate ``superstar'' incomes. Once workers reach retirement
age $J^{retired}$ income falls to a fixed level based on their wage
upon retirement. The income process is normalized to have a mean of
1, such that the price of housing $p^{h}_{t}$ reflects the house
price to income ratio.

Households may save or borrow liquid assets subject to a borrowing
constraint where each household can borrow against a proportion, $m$,
of the total value of the stock of housing they own

\begin{equation}
a_{t}\geq-mh^{own}_{t}p^{h}_{t}.
\end{equation}
Households may buy or sell housing at price, $p^{h}_{t},$ subject
to a minimum block size $\bar{h}$. All property sales, including
those that are sold as part of a bequest are subject to a transaction
cost $F_{sell}$ proportional to the total value of property sold,
which is reflected by the indicator function $I_{sell}$. If they
do not own any housing they must rent it from real estate firms at
price $p^{rent}_{t}$. The household's utility function is thus maximized
subject to the general budget constraint\footnote{Noting that households only interact with the rental market if they
own no housing stock of their own.}

\begin{equation}
a_{t+1}=\left(1+r^{a}_{t}\right)a_{t}+y_{t}-c_{t}-p^{rent}_{t}h^{rent}_{t}-\left(1-I^{sell}F^{sell}\right)p^{h}_{t}\varDelta h^{own}_{t}
\end{equation}
where $I^{sell}$ is an indicator variable for households that choose
to sell part or all of their current stock of housing.

We assume the households start with some initial wealth with the mean
of the distribution of initial liquid wealth, $a_{25}$, defined by
the average bequest of the household's at the end of their life.

\subsection{Household's Problem}

Each household's state vector is $\mathbf{s}_{t}=\left\{ a_{t},h^{own}_{t-1},z_{t}\right\} $,
where $a_{t}$ is the stock of liquid assets, $h^{own}_{t-1}$ is
the amount of housing they owned at the end of period $t-1$ and $z_{t}$
is the persistent component of their income. \rev{Let
$X_{t}\equiv\left(\Omega_{t},p^{h}_{t},p^{rent}_{t}\right)$ denote
the aggregate state in period $t$. In the transition experiment, households
take the future path $\{X_{\tau}\}_{\tau=t}^{T}$ as given; in the stationary
benchmark, the aggregate state is constant over time.} The value function for the
household is given by the maximum of the three possible housing tenure
options: rent (either because they owned no housing or sell it all
that period), or are owners who choose to adjust their stock of housing
or maintain their current level of housing. Households who choose
to rent (and thus sell whatever stock of housing they owned at the
beginning of the period) are able to freely adjust the amount of housing
they rent each period. The value function for a household that choose
to rent housing (and thus sell any previously owned housing stock)
at age age $j$ is

\begin{eqnarray}
\revmath{V^{rent}_{j}\left(\mathbf{s}_{t};X_{t}\right)} & = & \revmath{\underset{c_{t},a_{t+1},h^{rent}_{t}}{\mbox{max}}u\left(c_{t},h^{rent}_{t}\right)+\beta E\left[V_{j+1}\left(\mathbf{s}_{t+1};X_{t+1}\right)\mid z_{t},X_{t}\right]}\\
s.t\nonumber \\
a_{t+1} & = & \left(1+r^{a}_{t}\right)a_{t}+y_{t}-c_{t}-p^{rent}_{t}h^{rent}_{t}+\left(1-F^{sell}\right)p^{h}_{t}h^{own}_{t-1}\nonumber \\
a_{t+1} & \geq & 0\nonumber \\
h^{own}_{t} & = & 0\nonumber
\end{eqnarray}
The value function for home owners who choose not to adjust their
(non-zero) stock of housing in period $t$ is

\begin{eqnarray}
\revmath{V^{non-adjust}_{j}\left(\mathbf{s}_{t};X_{t}\right)} & = & \revmath{\underset{c_{t},a_{t+1}}{\mbox{max}}u\left(c_{t},h^{own}_{t}\right)+\beta E\left[V_{j+1}\left(\mathbf{s}_{t+1};X_{t+1}\right)\mid z_{t},X_{t}\right]}\\
s.t\nonumber \\
a_{t+1} & = & \left(1+r^{a}_{t}\right)a_{t}+y_{t}-c_{t}\nonumber \\
a_{t+1} & \geq & -mp^{h}_{t}h^{own}_{t}\nonumber \\
h^{own}_{t} & = & h^{own}_{t-1}\nonumber
\end{eqnarray}
finally, the value function for households who choose to buy or sell
housing in period $t$ is

\begin{eqnarray}
\revmath{V^{adjust}_{j}\left(\mathbf{s}_{t};X_{t}\right)} & = & \revmath{\underset{c_{t},a_{t+1},h^{own}_{t}}{\mbox{max}}u\left(c_{t},h^{own}_{t}\right)+\beta E\left[V_{j+1}\left(\mathbf{s}_{t+1};X_{t+1}\right)\mid z_{t},X_{t}\right]}\\
s.t\nonumber \\
a_{t+1} & = & \left(1+r^{a}_{t}\right)a_{t}+y_{t}-c_{t}-\left(1-I^{sell}F^{sell}\right)p^{h}_{t}\varDelta h^{own}_{t}\nonumber \\
a_{t+1} & \geq & -mp^{h}_{t}h^{own}_{t}\nonumber \\
h^{own}_{t} & = & h^{own}_{t-1}+\varDelta h^{own}_{t}\nonumber
\end{eqnarray}
The final value function is the maximum of each of the three choices
over housing tenure

\begin{equation}
\revmath{V_{j}\left(\mathbf{s}_{t};X_{t}\right)=\underset{}{\mbox{max}\left\{ V^{rent}_{j}\left(\mathbf{s}_{t};X_{t}\right),V^{adjust}_{j}\left(\mathbf{s}_{t};X_{t}\right),V^{non-adjust}_{j}\left(\mathbf{s}_{t};X_{t}\right)\right\} }}
\end{equation}
\rev{In the stationary benchmark, $X_{t+1}=X_{t}=X\left(\Omega\right)$,
so aggregate prices are constant and the expectation operator applies
only to next period's idiosyncratic income shock. Along a political
rational-expectations transition, the sequence $\{X_{t}\}_{t=0}^{T}$
is deterministic once the demographic shock is specified.}

The expected value function for the final year of life is the same
for all types of housing tenure and is given by the warm glow bequest
function given by equation 4 where total wealth is the sum of liquid
and illiquid wealth given the current housing price.

\begin{equation}
w=a+\left(1-F^{sell}\right)p^{h}_{t}h^{own}
\end{equation}


\subsection{Real Estate Firms}

Real estate firms borrow from households by issuing debt in the form
of the liquid asset to fund the purchase of housing, which is rented
out to households that do not own their own home. Real estate firms
operate in a perfectly competitive market, and thus the rental price
is pinned down in the quantitative implementation by the financing
wedge between borrowing and saving and a fixed cost of providing rental
services. We assume that real estate firms are able to borrow against
the full value of their housing stock. The rental price is therefore
given by

\begin{equation}
\revmath{p^{rent}_{t}=\phi^{rent}+r^{b}_{t}-r^{a}_{t},}
\end{equation}
where $\phi^{rent}$ is the per unit cost of providing rental services
and $r^{b}_{t}$ is the borrowing rate faced by landlords. \rev{In
the quantitative implementation, $r^{b}_{t}=r^{a}_{t}+\text{spread}$,
so rental prices are a derived object rather than an additional fixed-point
price.}

\subsection{Housing Supply}

The aggregate economy consists of unit measure of neighborhoods across
which all types of households are equally distributed. The neighborhood's
housing supply in period $t+1$ is determined in each neighborhood
by a simple majoritarian voting mechanism in which households vote
over the amount of permitted housing in an area. Only households who
currently reside in that neighborhood can vote to determine the next
periods level of new housing. Once the quantity of housing permits
has been determined by the households' collective vote within the
neighbourhood, the supply of housing is costlessly produced by a perfectly
competitive construction sector in $t+1$. This is similar to Kaplan
et al. (2020), who model a perfectly competitive construction sector
in which all rents from land permits accrue to the government.

An increase in the housing supply will reduce the price of housing
in the locality and households vote understanding that an increase
in housing permits will lead to a small persistent decrease in the
cost of housing. An increase in the housing supply will also affect
the price of rental properties by lowering the input price for the
real estate firms. As a result, each household in a neighborhood will
vote to increase (or decrease) the availability of housing permits
in their locality based on how a small, but permanent, change in house
prices will affect their expected welfare. Specifically, households
will vote to decrease the supply of housing if a small change in house
prices, $\epsilon^{p}$ , increases their expected lifetime utility
such that

\begin{equation}
\revmath{V_{j}\left(\mathbf{s};p^{h}+\epsilon^{p},\Omega\right)-V_{j}\left(\mathbf{s};p^{h},\Omega\right)+\sigma\geq0}
\end{equation}
where $\sigma$ represents a non-financial preference over the supply
of housing. A negative or positive $\sigma$ indicates a non-financial
preference for lower or higher housing supply, respectively. For example,
households may view new housing negatively as a result of congestion
effects on infrastructure, or positively because of the new amenities
that greater density could bring. Households may also hold ideological
views about housing supply, which this parameter would capture as
well.

In principle these preferences could vary by household, either due
to changes in age, productivity, or as a result of an idiosyncratic
shock. However, allowing for such heterogeneity would require introducing
a larger set of novel parameters that would then need to be calibrated.
To limit our degrees of freedom, we therefore calibrate these preferences
as a simple, fixed constant across all households regardless of age,
income or wealth. As a result, this parameter is effectively the average
non-financial preference for all households. We test this assumption
in Section 6 and find that the main results of the model do not substantively
change with alternative assumptions.

The impact of a small perturbation to the price of housing will vary
depending on household's ownership of liquid assets, housing and their
current level of productivity. If a small increase in the price of
housing increases their expected lifetime utility, then the left hand
side of equation 14 will be positive and thus they will vote for lower
supply of housing permits.

\subsection{The Medan Voter Theorem}

Applying the median voter theorem, equilibrium in the housing market
will occur when the median voter is indifferent to a change in the
price of housing.\footnote{While these preferences are in general well defined, it is important
to note that preferences over house prices are not monotonic. For
example, a household with a moderate amount of housing might prefer
a small increase in house prices as this will increase their expected
bequest, but would be better off if house prices fall substantially
as they would be able to consume more housing.} \rev{For a given demographic
state $\Omega$, the equilibrium house price $p^{h}\left(\Omega\right)$
is determined by the median voter equation}
\begin{equation}
\revmath{\sum_{j=25}^{80}\omega_{j}\int g_{j}\left(\mathbf{s};p^{h},\Omega\right)\Theta_{j}\left(\mathbf{s};p^{h},\Omega\right)d\mathbf{s}=0.5}
\end{equation}

and

\begin{equation}
\revmath{\Theta_{j}\left(\mathbf{s};p^{h},\Omega\right)=\mathbf{1}_{\left\{ V_{j}\left(\mathbf{s};p^{h}+\epsilon^{p},\Omega\right)-V_{j}\left(\mathbf{s};p^{h},\Omega\right)+\sigma\geq0\right\} }}
\end{equation}

\rev{in which $g_{j}\left(\mathbf{s};p^{h},\Omega\right)$ is the distribution
of age-$j$ households over idiosyncratic states, and $\Theta_{j}\left(\mathbf{s};p^{h},\Omega\right)$
indicates whether an age-$j$ household prefers higher house prices.
The median voter condition sets the population share voting for higher
house prices equal to 0.5.}

The change in expected lifetime utility from a small, persistent change
in the price of housing is determined by the households current portfolio
of assets and expected life-time income. Figure \ref{Figure 4} shows
how different household characteristics correlate with the decision
to vote for lower supply and thus higher house prices. Households
who currently rent or expect to have a high likelihood of purchasing
new housing in the near term will vote for housing supply to increase
thus decreasing the price of housing and relaxing their budget constraint.
These households tend to be younger, have low incomes and low wealth.
Looking at the distribution of voters according to the expected probability
that a given household will make a purchase of housing at some point
in their life, we can see that households that do not have a high
expected probability of purchasing housing over the rest of their
life-time will generally prefer a higher price of housing in order
to increase the size of their bequest and the resale value of their
stock of housing. By contrast the coalition of support for lower house
prices consists of two distinct groups of households: households that
anticipate buying additional housing (particularly in the near term)
and thus want to loosen their budget constraint. as well as households
that are currently renters because an increase in the supply of housing
also lowers the price of rental units.

The equilibrium price of housing will thus vary according to the relative
size of these groups of households. An exogenous increase in young
households, who disproportionately rent and thus prefer lower house
prices, will lead to an increase in the supply of housing. The price
of housing will thus fall until a sufficient number of renters become
homeowners, and accordingly switch their political support towards
higher house prices, thus making the median household indifferent
to a change in house prices.

In order to abstract from strategic voting across dynasties, we assume
different children and parents live in different neighborhoods and
thus do not effect the other's house market dynamics.\footnote{According to Fu (2019) 61.5\% adults live more than 10 miles away
from their mother.} Consequently even if an individual may benefit from a larger inheritance
if the price of their parents home increases, it will not affect their
vote on the optimal level of housing supply in their local neighborhood.\footnote{This assumption also makes our model robust to other forms of altruistic
giving. For example, in addition to receiving a warm glow from giving
to their descendants at death, households may care about their descendant's
welfare directly and thus may wish to sell housing in order to give
their children funds to purchase housing earlier. As families live
in different neighbourhoods, however, this would still not change
their voting behaviour. This is because parents with housing would
still wish to restrict housing supply in their neighbourhood in order
to receive the highest return on their house so they can maximise
their child's welfare.}

\rev{Households are forward-looking, and expectations enter the model
through the continuation values in their dynamic problem. In the main
transition experiment, households solve given a candidate deterministic
future path of house prices and demographics; political support is
then evaluated on the implied distributional path, and the equilibrium
transition is the path for which these objects are mutually consistent.
Alternative forecasting rules appear only in robustness exercises. We
also assume that households are unable to form coalitions when voting,
nor bargain political support across time periods and vote solely
based on their individual interest at time $t$. Once the quantity
of housing permits has been determined by the households' collective
vote the supply of housing is costlessly supplied by a perfectly competitive
construction sector. This is similar to Kaplan et al (2020) who assume
a perfectly competitive construction sector in which all rents from
land permits accrue to the government.}\footnote{This assumption allows us to abstract from any potential employment
benefits derived from additional demand for the construction sector. }

\begin{figure}[H]
\begin{centering}
\includegraphics[scale=0.08]{\string"Figures/4.1 VoteByYQ\string".eps}\includegraphics[scale=0.08]{\string"Figures/4.2 VoteByhQ\string".eps}
\par\end{centering}
\begin{centering}
\includegraphics[scale=0.08]{\string"Figures/4.3 VoteByAge\string".eps}\includegraphics[scale=0.08]{\string"Figures/4.4 DensityVote\string".eps}
\par\end{centering}
\noindent\raggedright{}\caption{Household Support for Higher House Prices}
\emph{\label{Figure 4}}
\end{figure}


\subsection{Equilibrium }

\rev{For ease of notation, let the distribution of household ages be
$\Omega=\{\omega_{25},...\omega_{80}\}$, where $\omega_{i}$ is the
share of the population at age $i$.}\footnote{We drop the $t$ subscript for simplicity.}
\rev{In the stationary benchmark, a political-economic equilibrium
indexed by $\Omega$ consists of age-specific value functions and policy
rules $\{V_{j},\chi_{j}\}_{j=25}^{80}$, age-specific state distributions
$\{g_{j}\}_{j=25}^{80}$, a house price $p^{h}\left(\Omega\right)$,
and a rental price $p^{rent}\left(\Omega\right)$ such that: (i) given
$X\left(\Omega\right)$, households solve the Bellman problems above;
(ii) the derived rent rule holds, $p^{rent}\left(\Omega\right)=\phi^{rent}+r^{b}-r^{a}$;
(iii) the voting condition $\sum_{j=25}^{80}\omega_{j}\int g_{j}\left(\mathbf{s};p^{h}\left(\Omega\right),\Omega\right)\Theta_{j}\left(\mathbf{s};p^{h}\left(\Omega\right),\Omega\right)d\mathbf{s}=0.5$
holds; and (iv) the law of motion for household states and household
formation leaves the cross-sectional distribution invariant and preserves
the age shares $\Omega$.}
\rev{In the transition experiment, the demographic path $\{\Omega_{t}\}_{t=0}^{T}$
is exogenous and the equilibrium object is a whole path $\{V_{j,t},\chi_{j,t},g_{j,t},p^{h}_{t},p^{rent}_{t}\}_{t=0}^{T}$
such that households optimize given the deterministic future aggregate
path, the age-specific distributions evolve consistently with the
induced policy rules and demographic law of motion, the derived rent
rule holds period by period, and the voting condition evaluated on
the implied distributional path is satisfied in every period. The
full numerical computational method is outlined
in detail in the Appendix.}

\section{Calibration}

We calibrate the model to capture salient features of the US housing
market in the early 2000s prior to the housing boom and bust. The
model period is one year, households work for 41 periods (age 25 to
65) and die after 56 periods (age 80). The risk aversion parameter
is set to 2, as is standard in the macroeconomics literature. The
income process consists of the parameters for the deterministic age-profile,
the persistent AR(1) component, and the transitory i.i.d. component
of income. We estimate the parameters of the deterministic and stochastic
20 income processes using data from the Panel Study of Income Dynamics.

The retirement income replacement rate is set at 50\% of the final
period's non-transitory income, based on Díaz et al. (2008). The risk-free
interest rate is set at 2\%, while the mortgage interest rate is 5\%.
The maximum loan-to-value (LTV) is set at 0.9 in line with Greenwald
(2018). The spread between the saving rate and the mortgage rate is
set at 3\% matching the 30 year average difference between the two.

{\scriptsize{}
\begin{table}[H]
{\scriptsize\caption{Parameter Values}
}{\scriptsize\par}
\centering{}{\scriptsize{}%
\begin{tabular}{cc>{\centering}p{3cm}>{\centering}p{3cm}}
\toprule
{\scriptsize Parameter} & {\scriptsize Value} & {\scriptsize Interpretation} & {\scriptsize Source}\tabularnewline
\midrule
{\scriptsize$J$} & {\scriptsize 56} & {\scriptsize Periods of life} & {\scriptsize KMV (2020)}\tabularnewline
{\scriptsize$J_{ret}$} & {\scriptsize 41} & {\scriptsize Period of working life} & {\scriptsize KMV (2020)}\tabularnewline
{\scriptsize$r^{a}$} & {\scriptsize 0.02} & {\scriptsize Risk free rate} & {\scriptsize 30 year average}\tabularnewline
{\scriptsize$\beta$} & {\scriptsize$0.978$} & {\scriptsize Discount rate} & {\scriptsize Internally Calibrated}\tabularnewline
{\scriptsize$\gamma$} & {\scriptsize 2} & {\scriptsize Risk aversion} & {\scriptsize KMV (2020)}\tabularnewline
{\scriptsize$\zeta$} & {\scriptsize 0.75} & {\scriptsize Housing weight in } & {\scriptsize Internally Calibrated}\tabularnewline
{\scriptsize$F^{sell}$} & {\scriptsize 0.07} & {\scriptsize Cost of selling house} & {\scriptsize KMV (2012020)}\tabularnewline
{\scriptsize$m$} & {\scriptsize 0.9} & {\scriptsize Maximum LTV} & {\scriptsize Guerrieri and Iacoviello (2017)}\tabularnewline
{\scriptsize$\psi$} & {\scriptsize 10} & {\scriptsize Bequest scaling} & {\scriptsize Internally Calibrated}\tabularnewline
{\scriptsize$\overline{w}$} & {\scriptsize$150$} & {\scriptsize Bequest parameter} & {\scriptsize Internally Calibrated}\tabularnewline
{\scriptsize$\sigma$} & {\scriptsize$-2.4$} & {\scriptsize Pro-supply Non-financial preference} & {\scriptsize Internally Calibrated, shifts 15\% of the voters in steady
state}\tabularnewline
{\scriptsize$\phi^{rent}$} & {\scriptsize$0.05$} & {\scriptsize Rental per unit cost} & {\scriptsize Internally Calibrated}\tabularnewline
{\scriptsize$\theta^{rent}$} & {\scriptsize 0.9} & {\scriptsize Rental discount} & {\scriptsize Mian et al (2017)}\tabularnewline
\bottomrule
\end{tabular}}{\scriptsize\par}
\end{table}
}{\scriptsize\par}

We set the discount rate, rental discount, utility housing weight
and real estate markup $\left\{ \beta,\sigma,\zeta,\overline{w},\psi,\phi^{rent}\right\} $
using simulated method moments to match the following moments.

\begin{table}[H]
\caption{Calibration}

\centering{}%
\begin{tabular}{c>{\centering}p{2cm}>{\centering}p{2cm}}
\toprule
 & Model & Data\tabularnewline
\midrule
\midrule
Share of owners with a mortgage & 0.55 & 0.57\tabularnewline
Share of renters & 0.38 & 0.38\tabularnewline
House price to Income Ratio & 6.0 & 6.0\tabularnewline
Home-ownership rate 25-35 & 0.38 & 0.40\tabularnewline
Home-ownership rate 45-55 & 0.72 & 0.70\tabularnewline
Home-ownership rate 65-75 & 0.71 & 0.75\tabularnewline
\bottomrule
\end{tabular}
\end{table}

The price point at which the median voter is indifferent to additional
housing supply will be affected by $\sigma$, the households ideological
preference over the supply of housing, this allows households to have
non-financial preferences over the supply of housing. We calibrate
$\sigma$ such that the housing market is in equilibrium in 2000.
Our calibration requires a non-financial preference in favor of approving
more homes as homeowners almost always make up a large majority of
the voting population.

\section{Quantitative Analysis}

\subsection{Steady State}

\rev{We begin with a stationary benchmark. In this section, ``steady
state'' refers to a sequence of stationary political-economic equilibria
indexed by the age distribution of households, rather than to a transition
path with changing demographic shares. For each year's demographic
composition, we hold non-demographic parameters fixed, choose household
formation so that the age structure is constant in the absence of
further shocks, and solve for the house price that clears the political
housing-supply condition. Equivalently, this section traces the mapping
from $\Omega$ to the associated stationary equilibrium price $p^{h}\left(\Omega\right)$.}
\rev{We calculate this stationary benchmark using demographic data
from the US over the period 1950-2020.}
During this time there was substantial variation in the demographic
structure - driven by the generation born in the post-war spike and
how it reverberated over time - as shown in figure \ref{Figure 3}
at the beginning of the paper. For example, the mean age starts at
30 in 1950 and eventually falls as those born during the baby boom
increases the share of young people in the labor and housing market.
The mean age in the labor market troughs at 28 in 1970 as the baby
boom subsides. The mean age then continues to steadily rise to 39
in 2020 as the those born during the baby boom ages.

\rev{This benchmark is useful because it isolates the comparative-static
effect of demographic composition on the median voter's preferred housing
supply. The dynamic transition analysis below asks a different question:
how a temporary baby boom propagates when households form model-consistent
expectations about future prices and future political support.}

Since homeownership rates increases with age, we find that when the
median population age increases the aggregate homeownership rate rises.
This increases political support for a reduction in the housing supply
which increases house prices. House prices rise until their is enough
political support for additional housing supply (via an increase in
households who rent) such that the median voter is again indifferent
to a small change in house prices.

Figure \ref{Figure 5} shows the model's political-economic equilibrium
cost of housing over the post war period along side data on real rents.\footnote{In our model's environment real rents move inline with the cost of
housing, while abstracting from the impact of interest rates which
will impact on the empirical price of housing.} This equilibrium cost represents the effects of demographic shocks
on the coalition of voters only on the housing market, and thus ignores
other factors such as changes in the business cycle, technology and
financial crises. Accordingly one should not expect it to closely
track actual costs that are affected by these non-demographic factors.
Despite this the equilibrium results manage fit the post war trend
in real rents quite closely.\footnote{We use rents instead of house prices as a measure of housing costs
to avoid the measure being distored by the secular decline in interest
rates over this period. Real rents are calculated using the average
rent of a primary residence across US cities divided by headline CPI.}In the model the post war spike in the birth rate leads to a significant
fall in the equilibrium house price starting in the 1970s as these
households, who are mostly young renters, enter the housing market.
As these households age and become homeowners they shift towards opposing
the construction of new housing to protect the value of their stock
of housing. Real rents show a similar trend peaking slightly earlier
in the 1950s, before trending downwards for the next three decades
reaching a low point in the 1980s then increasing steadily up until
2020.
\begin{figure}[H]
\begin{centering}
\includegraphics[scale=0.1]{\string"Figures/5b. PriceIncomeRatio with Real Rents\string".eps}
\par\end{centering}
\noindent\raggedright{}\caption{Steady State Equilibrium House Price\emph{ }}
\emph{\label{Figure 5}}
\end{figure}

Using demographic projections out to 2100 we can estimate the impact
the on-going aging of the US economy will have on the political economy
of the US housing market as shown in figure \ref{Figure 6}. The IMF
calculate demographic projections for the US population subject to
three different assumptions about the level of immigration, of primarily
young households, over the next 80 years. We find that for low to
medium assumptions about immigration the aging of the US population
will continue to increase the relative size of the electorate who
support higher house prices - at least at the initial level of house
prices. As a result, the equilibrium level of house prices increases
over the next 80 years. For high level assumptions about immigration
we find results similar to that of a spike in the birth rate, this
is because immigrants tend to be younger and would change the demographic
structure towards a lower median-age which would initially result
in more political support for increase in the housing supply. The
high immigration projection includes some long run cyclical dynamics
as young immigrants eventually age and swing between the NIMBY and
anti-NIMBY coalitions.
\rev{We treat Figures \ref{Figure 5} and \ref{Figure 6} as stationary
benchmark objects. They discipline the long-run demographic mechanism,
but they are not substitutes for the fully dynamic political rational-expectations
transition studied next.}

\begin{figure}[H]
\begin{centering}
\includegraphics[scale=0.1]{\string"Figures/6. Price Income Ratio Scenario\string".eps}
\par\end{centering}
\noindent\raggedright{}\caption{Steady State Equilibrium House Price Projections}
\emph{\label{Figure 6}}
\end{figure}


\subsection{Impulse Response Functions}

\rev{This subsection studies the transition generated by a temporary
baby boom when households form political rational expectations over
the future path of house prices and housing supply. The dynamic
transition complements the stationary benchmark by tracing how the
boom cohort reshapes housing demand, political support for additional
supply, and the distribution of homeownership and wealth over time.}

\rev{The transition is organized around three objects:
the path of the boom cohort through renting and homeownership, the
induced evolution of political support for additional housing, and
the resulting path of house prices, homeownership, and wealth across
cohorts. If reduced-form expectation exercises are retained, the relevant
comparison is between the political rational-expectations transition
and those alternative transition closures, not between different stationary
benchmarks.}
\rev{Computationally, we proceed in four steps. First, we guess a finite
path of house prices. Second, given that path and the exogenous demographic
path, we solve the household problem backward. Third, we simulate the
cross-sectional distribution forward and evaluate political support
period by period using the implied perturbed-price value objects. Finally,
we update the guessed price path and iterate until the price path, the
distributional path, and the political residuals are mutually consistent.}
\begin{figure}[H]
\begin{centering}
\includegraphics[scale=0.08]{\string"Figures/7.1 Bith\string".eps}\includegraphics[scale=0.08]{\string"Figures/7.2 Average Age\string".eps}
\par\end{centering}
\begin{centering}
\includegraphics[scale=0.08]{\string"Figures/7.3 Smoothed House Price\string".eps}\includegraphics[scale=0.08]{\string"Figures/7.4 Home30\string".eps}
\par\end{centering}
\noindent\raggedright{}\caption{Impulse Response Functions for a Temporary Birthrate Increase.
\rev{[Redo from political RE transition.]}}
\emph{\label{Figure 7}}
\end{figure}

\rev{[Figure \ref{Figure 8}: redo from political RE transition.]}

\begin{figure}[H]
\begin{centering}
\includegraphics[scale=0.08]{\string"Figures/8.1 Generations 1\string".eps}\includegraphics[scale=0.08]{\string"Figures/8.2 Generations 3\string".eps}
\par\end{centering}
\begin{centering}
\includegraphics[scale=0.08]{\string"Figures/8.3 NW2\string".eps}\includegraphics[scale=0.08]{\string"Figures/8.4 NW1\string".eps}
\par\end{centering}
\noindent\raggedright{}\caption{Impulse Response Functions for a Temporary Birthrate Increase.
\rev{[Redo from political RE transition.]}}
\emph{\label{Figure 8}}
\end{figure}


\subsection{Permanent Demographic Shock}

We next consider the impact of a permanent demographic shock as an
alternative to a transitory shock. In this case the formation rate
of new households falls 10 per cent forever. While the fall in the
birth rate is permanent the effect on the model is temporary: once
the new formation rate has time to filter through the demographic
structure of the economy, the relative sizes of the different generations
will return to the original steady state of the model. We find that
this permanent reduction in the formation of new households causes
a transitory increase in the average age that lasts 56 years. This
increase in the average age increases the relative size of the older
generations. Since these older generations are predominantly home
owners and more likely to prefer a lower level of housing supply,
we see an increase in the equilibrium house price and a corresponding
reduction in the homeownership rate and the consumption of housing.
\rev{[Figure \ref{Figure 9}: redo from political RE transition if retained.]}
\begin{figure}[H]
\begin{centering}
\includegraphics[scale=0.08]{\string"Figures/9.1 Average Age with Negative Permanent Shock\string".eps}\includegraphics[scale=0.08]{\string"Figures/9.2 Smoothed House Prices with Negative Permanent Shock\string".eps}
\par\end{centering}
\begin{centering}
\includegraphics[scale=0.08]{\string"Figures/10.1 Generations1\string".eps}\includegraphics[scale=0.08]{\string"Figures/10.2 Generations2\string".eps}
\par\end{centering}
\noindent\raggedright{}\caption{Impulse Response Function for Permanent Birthrate Fall.
\rev{[Redo from political RE transition if retained.]}}
\emph{\label{Figure 9}}
\end{figure}

\textcompwordmark{}

\section{Robustness}

\rev{The robustness exercises serve two distinct purposes. Some are
stationary-benchmark checks that ask whether the mapping from demographics
to the stationary political equilibrium is sensitive to parameter choices
or voting weights. Others, if retained, are alternative transition closures.
In that latter case the natural comparison is with the political rational-expectations
transition itself rather than with the stationary benchmark.}

We first consider how these results are affected by choice of calibration
of the model. Varying the size of non-financial preferences over housing
supply largely does not affect the dynamics of the model. While a
stronger (weaker) non-financial preference for housing supply does
lower (increase) the steady state house price to income ratio the
impacts on the dynamics of the model are small. Regardless of whether
the non-financial preferences are increased or decreased the impact
of a temporary baby boom on house prices is broadly unchanged. This
is because while the mean level of support changes under the alternative
calibration differs the marginal impact of a boom generation is mostly
unchanged.
\rev{[Figure \ref{Figure 11}: redo from political RE transition if retained.]}

\begin{figure}[H]
\begin{centering}
\includegraphics[scale=0.1]{\string"Figures/11. Smoothed House Price Params\string".eps}
\par\end{centering}
\noindent\raggedright{}\caption{Steady State Equilibrium House Price.
\rev{[Redo from political RE transition if retained.]}}
\emph{\label{Figure 11}}
\end{figure}


\subsection{Alternative Voting Mechanisms}

\rev{This subsection returns to stationary benchmark objects. The alternative
voting rules change the mapping from demographics to the stationary
equilibrium house price, rather than introducing a different transition-expectations
structure.}

Finally we consider how variations in the voting mechanism may affect
the results for the equilibrium house prices. We consider two alternatives:
weighting household voting power by turnout rates and weighting households
political power by the log of their net worth.

Weighting households vote by their turnout rate (in US elections)
by age does not materially change the results. Younger households
tend to have a lower voting rates, a result which has been constant
across time. However, the overall estimated path of house prices does
not change.

Alternatively, weighting by household's political power by using the
log of their net worth, has a larger impact on the path of projected
house prices. The fall in house prices that is otherwise estimated
to have occurred in the 1970s, when the baby boom generation are mostly
renters occurs later. Although they are more numerous than previous
generations, their relatively low level of net worth means they a
smaller degree of political leverage compared with the base case.
As the economy ages, the increase in the price to income ratio is
even larger as the aging population transition towards both home ownership
and greater political power as their net worth increases. This alternative
voting mechanism does not qualitatively change the overall results,
but it does make it quantitatively larger in the 21st century.

\begin{figure}[H]
\begin{centering}
\includegraphics[scale=0.1]{\string"Figures/12. PriceIncome Ratio Alternatives\string".eps}
\par\end{centering}
\noindent\raggedright{}\caption{Steady State Equilibrium House Price}
\emph{\label{Figure 12}}
\end{figure}


\section{Conclusion}

\rev{Local planning regulations and the political-economic constraints
that bind them have long been defining features of housing markets
across high-income countries. Given the importance of housing in homeowner's
financial portfolios, it is natural that some households might seek
to protect this value by voting in favor of regulations that restrict
the supply of new housing. In our stationary benchmark, post-war demographic
changes shift the age composition of the electorate, change political
support for additional housing supply, and thereby move the stationary
equilibrium house-price-to-income ratio. The benchmark implies a sizable
demographic component in post-war house-price movements and points
to further upward pressure as the population ages. The natural dynamic
object for studying temporary demographic shocks is the political rational-expectations
transition, which traces how housing demand, voting incentives, and
prices co-evolve along the full path.}

Unpacking the distributional impacts of this generational shift, we
find that a one-off increase in the birth rate generally increases
the relative consumption of the boom generation as they use their
disproportionate voting power to improve their financial position
throughout their life cycle. However, this comes at the cost of other
generations: the consumption of the parent's generation initially
falls during the boom, and the net worth of the baby boomers children
is lower over their entire life cycle. This highlights how demographic
shocks can have long-lasting effects, even after the baby boom generation
has completely died off. Eventually, the housing market returns to
equilibrium, but the effects of a baby boom can distort the housing
market for both themselves and future generations.

\section*{References}

Ahlfeldt, G. M. (2011). Blessing or curse? Appreciation, amenities
and resistance to urban renewal. Regional Science and Urban Economics,
41(1):32\textendash 45.

Ahlfeldt, G. M. and Maennig, W. (2015). Homevoters vs. leasevoters:
A spatial analysis of airport effects. Journal of Urban Economics,
87:85\textendash 99

Albouy, D., Behrens, K., Robert-Nicoud, F., \& Seegert, N. (2019).
The optimal distribution of population across cities. Journal of Urban
Economics, 110, 102-113.

Anagol, S., Ferreira, F. V., \& Rexer, J. M. (2021). Estimating the
economic value of zoning reform (No. w29440). National Bureau of Economic
Research.

Attanasio, O. P., Bottazzi, R., Low, H. W., Nesheim, L., \& Wakefield,
M. (2012). Modelling the demand for housing over the life cycle. Review
of Economic Dynamics, 15(1), 1-18.

Alesina, A., \& Rodrik, D. (1994). Distributive politics and economic
growth. The quarterly journal of economics, 109(2), 465-490.

Bunten, D. (2017). Is the rent too high? aggregate implications of
local land-use regulation.

Bosetti, N., and Sims, S. (2016). \textquotedbl Stopped: Why people
oppose development in their back yard\textquotedbl{} Centre for London

Chatterjee, S., \& Eyigungor, B. (2017). A tractable city model for
aggregative analysis. International economic review, 58(1), 127-155.

Cheung, R. and Meltzer, R. (2013). Homeowners associations and the
demand for local land use regulation. Journal of Regional Science,
53(3):511\textendash 534.

Dehring, C. A., Depken II, C. A., and Ward, M. R. (2008). A direct
test of the homevoter hypothesis. Journal of Urban Economics, 64(1):155\textendash 170.

Duranton, G., \& Kerr, W. R. (2015). The logic of agglomeration (No.
w21452). National Bureau of Economic Research.

Duranton, G., \& Puga, D. (2023). Urban growth and its aggregate implications.
Econometrica, 91(6), 2219-2259.

Eggleston, J., Hays, D., Munk, R., \& Sullivan, B. (2020). The Wealth
of Households, 2017. Washington, DC, USA: US Department of Commerce,
US Census Bureau.

Einstein, K. L., Palmer, M., \& Glick, D. M. (2019). Who participates
in local government? Evidence from meeting minutes. Perspectives on
politics, 17(1), 28-46.

Fernandez-Villaverde, J., \& Krueger, D. (2011). Consumption and saving
over the life cycle: How important are consumer durables?. Macroeconomic
dynamics, 15(5), 725-770.

Fischel, W. A. (1999). Does the American way of zoning cause the suburbs
of metropolitan areas to be too spread out?. Governance and opportunity
in metropolitan America, 15, 1-19.

Fischel, W. A. (2002). The homevoter hypothesis: How home values influence
local government taxation, school finance, and land-use policies.
Cambridge, MA: Harvard University Press.

Frieden, B. J. (1979). Environmental protection hustle. United States.

Fu, C. H. (2019). Living arrangement and caregiving expectation: The
effect of residential proximity on inter vivos transfer. Journal of
Population Economics, 32(1), 247-275.

Duranton, G. \& Kerr, W. R. (2015). The Logic of Agglomeration\textquotedbl .
NBER Working Papers 21452, National Bureau of Economic Research, Inc.

Glaeser, E. L., Gyourko, J., \& Saks, R. E. (2006). Urban growth and
housing supply. Journal of economic geography, 6(1), 71-89.

Greenwald, D. (2018). The mortgage credit channel of macroeconomic
transmission. Research Paper no. 5184-16, Sloan School Management,
Massachusetts Inst. Technology.

Greenwood, J., Seshadri, A., \& Vandenbroucke, G. (2005). The baby
boom and baby bust. American Economic Review, 95(1), 183-207.

Gyourko, J., Mayer, C., \& Sinai, T. (2013). Superstar cities. American
Economic Journal: Economic Policy, 5(4), 167-199.

Hankinson, M. (2018). When do renters behave like homeowners? High
rent, price anxiety, and NIMBYism. American political science review,
112(3), 473-493.

Herkenhoff, K. F., Ohanian, L. E., \& Prescott, E. C. (2018). Tarnishing
the golden and empire states: Land-use restrictions and the US economic
slowdown. Journal of Monetary Economics, 93, 89-109.

Hilber, C. A. and Robert-Nicoud, F. (2013). On the origins of land
use regulations: Theory and evidence from us metro areas. Journal
of Urban Economics, 75:29\textendash 43.

Iacoviello, M. (2005). House prices, borrowing constraints, and monetary
policy in the business cycle. American economic review, 95(3), 739-764.

Iacoviello, M., \& Pavan, M. (2013). Housing and debt over the life
cycle and over the business cycle. Journal of Monetary Economics,
60(2), 221-238.

Jiang, B. (2018). Homeownership and voter turnout in us local elections.
Journal of Housing Economics, 41, 168-183.

Justiniano, A., Primiceri, G. E., \& Tambalotti, A. (2015). Household
leveraging and deleveraging. Review of Economic Dynamics, 18(1), 3-20.

Kaplan, G., Mitman, K., \& Violante, G. L. (2020). The housing boom
and bust: Model meets evidence. Journal of Political Economy, 128(9),
3285-3345.

Kahn, M. E., Vaughn, R., \& Zasloff, J. (2010). The housing market
effects of discrete land use regulations: Evidence from the California
coastal boundary zone. Journal of Housing Economics, 19(4), 269-279.

Kindermann, F., \& Krueger, D. (2022). High marginal tax rates on
the top 1 percent? Lessons from a life-cycle model with idiosyncratic
income risk. American Economic Journal: Macroeconomics, 14(2), 319-366.

Krusell, P., \& Smith, A. A., Jr. (1998). Income and wealth heterogeneity
in the macroeconomy. Journal of Political Economy, 106(5), 867\textendash 896.

Kulka, A., Sood, A., \& Chiumenti, N. (2022). How to increase housing
affordability: Understanding local deterrents to building multifamily
housing (No. 22-10). Working Papers.

Li, X. (2022). Do new housing units in your backyard raise your rents?.
Journal of Economic Geography, 22(6), 1309-1352.

Mankiw, N. G., \& Weil, D. N. (1989). The baby boom, the baby bust,
and the housing market. Regional science and urban economics, 19(2),
235-258.

Mast, E. (2024). Warding off development: Local control, housing supply,
and nimbys. Review of Economics and Statistics, 106(3), 671-680.

Manville, M., Monkkonen, P., \& Lens, M. (2020). It\textquoteright s
time to end single-family zoning. Journal of the American Planning
Association, 86(1), 106-112.

Mumtaz, H. and ¦ustek, R., 2023. Global house prices since 1950. Working
Paper.

Oliver, J. E., \& Ha, S. E. (2007). Vote choice in suburban elections.
American Political Science Review, 101(3), 393-408.

Ortalo-Magné, F., \& Prat, A. (2014). On the political economy of
urban growth: Homeownership versus affordability. American Economic
Journal: Microeconomics, 6(1), 154-181

Parkhomenko, A. (2023). Local causes and aggregate implications of
land use regulation. Journal of Urban Economics, 138, 103605.

Persson, T., \& Tabellini, G. (1994). Does centralization increase
the size of government?. European Economic Review, 38(3-4), 765-773.

Pennington, K. (2021). Does building new housing cause displacement.
The supply and demand effects of construction in san francisco.

Sahn, A. (2024). Racial diversity and exclusionary zoning: Evidence
from the great migration. Journal of Politics.

Scally, C. P., \& Tighe, J. R. (2015). Democracy in action?: NIMBY
as impediment to equitable affordable housing siting. Housing Studies,
30(5), 749-769.

Song, J. (2021). The effects of residential zoning in US housing markets.
Available at SSRN 3996483.

Titman, S., \& Zhu, G. (2024). City characteristics, land prices and
volatility. Journal of Urban Economics, 140, 103645.

Trounstine, J. (2020). The geography of inequality: How land use regulation
produces segregation. American Political Science Review, 114(2), 443-455.

Turner, M. A., Haughwout, A., \& Van Der Klaauw, W. (2014). Land use
regulation and welfare. Econometrica, 82(4), 1341-1403.

Whittemore, A. H., \& BenDor, T. K. (2018). Talking about density:
An empirical investigation of framing. Land use policy, 72, 181-191.

Yang, F. (2009). Consumption over the life cycle: How different is
housing?. Review of Economic Dynamics, 12(3), 423-443.

\newpage{}

\section*{Appendix}

The numerical solution for our model follows these steps.
\begin{enumerate}
\item The household's value functions are solve recursively for a fixed
value of $p^{h}$ starting at tin the final period of life and working
backwards using equations (8) - (11). This provides solutions for
the household value functions $\left\{ V^{rent}_{j},V^{adjust}_{j},V^{non-adjust}_{j}\right\} $
and the resulting decision rules for consumption, housing and liquid
assets as a function of the households idiosyncratic state variables.
\item The rental price of housing is pinned down by the cost of housing
$p^{h}$.
\item We then calculate how these value functions change if we permanently
increase the price of housing by a small amount, $\epsilon$. The
difference in the value functions is used to calculate the preferences
over housing supply across households, $\Theta$.
\item Next we solve for the density of households across the life cycle
using this decision rules to solve for the distribution of households
across the idiosyncratic state variables.
\item We then use these solutions to calculate the net-support for housing
supply which is defined the integrating the density of households
over their preferences for housing supply. If the net support is close
to $0.5$ we stop here, if the net support for housing supply is higher
(lower) than $0.5$ we lower (increase) the assumed price of housing
in step 1 and repeat the process until the median voter equation is
satisfied.
\item Once the housing market has been ``cleared'' by the political process,
we iterate forward one period updating the age of each household and
repeat the steps 1 through 5.
\end{enumerate}

\end{document}
